\documentclass[12pt,a4paper]{report}
\usepackage[utf8]{inputenc}
\usepackage[T1]{fontenc}
\usepackage[spanish]{babel}
\usepackage{csquotes}
\usepackage{geometry}
\usepackage{xcolor}
\usepackage{microtype}
\usepackage[colorlinks=true,urlcolor=blue,linkcolor=black,citecolor=black]{hyperref}
\geometry{margin=2.5cm}
\renewcommand{\thesection}{\arabic{section}}
\renewcommand{\thesubsection}{\arabic{section}.\arabic{subsection}}
\begin{document}
\sloppy

\begin{center}{\LARGE \textbf{Dictamen de Evaluación de Tesis}} \\ [0.5cm]\end{center}

\noindent\begin{tabular*}{\textwidth}{@{}l@{\extracolsep{\fill}}l@{}}
\textbf{Evaluador:} & {{EVALUADOR}} \\
\textbf{Rol:} & {{ROL}} \\
\textbf{Institución:} & {{INSTITUCION}} \\
\textbf{Fecha de Emisión:} & {{FECHA}} \\
\end{tabular*}

\tableofcontents

\hrule\vspace{0.5cm}

\clearpage

\section{Alcance de la Evaluación}

\noindent {{TEXTO_INTRO_IA}}
\vspace{0.5cm}

A continuación se detalla la evaluación sección por sección.


\section{Cumplimiento de Normativa}

El propósito de esta sección es verificar el cumplimiento integral de las obligaciones normativas derivadas del corpus legal peruano vigente al año 2026 en materia de formulación de tesis e investigación. La base normativa está al amparo de:

{{OBSERVACIONES_AQUI}}


\section{Respuesta a las Necesidades del Planeta}

Los Objetivos de Desarrollo Sostenible (ODS), adoptados por la ONU en 2015, son un llamado global para erradicar la pobreza, proteger el planeta y garantizar paz y prosperidad para todas las personas antes de 2030. Los 17 ODS son interdependientes: las acciones en un ámbito afectan a los demás y el desarrollo debe equilibrar la sostenibilidad social, económica y ambiental. Esta rúbrica permite evaluar la contribución potencial de un proyecto de tesis según los cinco ejes de la Agenda 2030.: \textbf{Personas}, \textbf{Planeta}, \textbf{Prosperidad}, \textbf{Paz} y \textbf{Alianzas}. Los 17 Objetivos de Desarrollo Sostenible, el agrupamiento por ejes de la Agenda 2030 y la escala de valoración deben incorporarse en esta sección como tablas de apoyo antes de la versión final del documento. La referencia normativa está con base en: Naciones Unidas, \textit{Transformar nuestro mundo: la Agenda 2030 para el Desarrollo Sostenible}, Resolución A/RES/70/1, 25 de septiembre de 2015. PNUD \url{https://www.undp.org/sustainable-development-goals}.

{{OBSERVACIONES_AQUI}}


\section{Respuesta a las Necesidades del País}

El nivel de complejidad metodológica del proyecto de tesis debe alinearse con la ley universitaria vigente y responder a las necesidades del país. El Plan Estratégico de Desarrollo Nacional (PEDN) al 2050, aprobado por el CEPLAN, órgano rector del SINAPLAN, es el principal instrumento de planeamiento estratégico del Perú. El PEDN sitúa a las personas y su vida en comunidad en el centro de su enfoque y define una estrategia de mediano y largo plazo para un desarrollo armónico y sostenible. Se organiza en cuatro Objetivos Nacionales (ON) —Desarrollo de las personas, Territorio sostenible, Competitividad e innovación, y Democracia y paz— que expresan cambios en el bienestar de la población bajo un enfoque sistémico. Esta rúbrica permite al dictaminante evaluar la contribución potencial de un proyecto de tesis a estos Objetivos Nacionales, a sus objetivos específicos y a los retos estructurales del PEDN, articulados con la Visión del Perú al 2050 y las 35 Políticas de Estado del Acuerdo Nacional. La Referencia normativa considerada es: Centro Nacional de Planeamiento Estratégico (CEPLAN), \textit{Plan Estratégico de Desarrollo Nacional al 2050}, aprobado por D.\,S.\,N.\textsuperscript{o}~095-2022-PCM, actualizado en noviembre de 2024. Visión del Perú al 2050, aprobada por consenso en el Foro del Acuerdo Nacional, abril de 2019. Fuente: \url{https://peru2050.ceplan.gob.pe}. En esta sección deben incorporarse los Objetivos Nacionales del PEDN al 2050, los ejes que estructuran la Visión del Perú al 2050 y la escala de valoración correspondiente como material de apoyo para el análisis.


{{OBSERVACIONES_AQUI}}


\section{Ajuste Epistemológico del Proyecto de Tesis}

El ajuste epistemológico de un proyecto de tesis exige que exista una alineación lógica y verificable entre los supuestos filosóficos del investigador, los objetivos de la investigación y los métodos seleccionados para alcanzarlos. La epistemología constituye un fundamento científico indispensable que orienta la adquisición, validación y aplicación del conocimiento en la investigación de la Ingeniería Civil. Un proyecto epistemológicamente coherente asegura que los compromisos ontológicos sobre la naturaleza de los fenómenos investigados, la orientación epistemológica sobre la validez del conocimiento y las implementaciones metodológicas se articulen de forma integrada y justificable. La presente rúbrica permite al dictaminante evaluar si el proyecto de tesis presenta coherencia y justificación filosófico-metodológica, flexibilidad e innovación, identificación de controversias y limitaciones, y comunicación científica mejorada en el contexto de las especialidades de la Ingeniería Civil: estabilización de suelos, pavimentos, cimentaciones, mecánica de materiales, investigación de nuevos materiales, entre otras. Las dimensiones de evaluación epistemológica y su escala de valoración deben incorporarse en esta sección como tablas de apoyo. Se toma como referencia de fundamentación el Capítulo 2 del libro \enquote{Metodología de la Investigación para la Estabilización Geotécnica de Suelos} (Arbulú \& Solorzano, 2026), con énfasis en la sección \textit{Importancia del Ajuste Epistemológico en Investigación Geotécnica} (subsecciones: Coherencia y Justificación, Flexibilidad e Innovación, Identificación de Controversias y Limitaciones, y Comunicación Mejorada). 


{{OBSERVACIONES_AQUI}}


\section{Evaluación de la Formulación del Problema de Investigación}

La problematización es el proceso intelectual mediante el cual el tesista identifica, delimita, contextualiza y estructura un problema de investigación, distinguiéndolo de uno de aplicación profesional y dándole sustento empírico y lógico previo a la hipótesis, la justificación y las preguntas de investigación. En Ingeniería Civil, exige rigor técnico en la descripción de fenómenos geotécnicos, estructurales, viales o de materiales, y sigue la secuencia lógica \textbf{Diagnóstico}~$\rightarrow$~\textbf{Pronóstico}~$\rightarrow$~\textbf{Control del Pronóstico}. El \textit{diagnóstico} describe la situación problemática actual con evidencia verificable; el \textit{pronóstico} proyecta las consecuencias de no intervenir; y el \textit{control del pronóstico} anticipa la estrategia de investigación para abordar el problema. Esta rúbrica permite evaluar la calidad, consistencia y completitud de la problematización, verificando que el planteamiento sea técnicamente preciso, no ambiguo, debidamente sustentado y lógicamente articulado como antecedente de la hipótesis, la justificación y la formulación del problema. La secuencia Diagnóstico--Pronóstico--Control del Pronóstico es una estructura lógico-argumental habitual en la tradición académica latinoamericana para formular el planteamiento del problema en proyectos de investigación de pregrado en Ingeniería Civil; su evaluación mediante esta rúbrica se centra en la consistencia interna, la precisión técnica y la suficiencia argumental como antecedentes de la hipótesis, la justificación y la formulación del problema de investigación.\par



{{OBSERVACIONES_AQUI}}


\section{Evaluación de las Justificaciones}

La justificación de un proyecto de tesis constituye el componente argumentativo que demuestra la pertinencia, la necesidad y el valor de la investigación propuesta. Una justificación rigurosa trasciende la simple declaración de intenciones y articula, con evidencia y razonamiento lógico, por qué la investigación merece ser ejecutada, financiada y difundida. La presente rúbrica permite al dictaminante evaluar las cinco dimensiones de la justificación que todo proyecto de investigación debe abordar de manera diferenciada: la \textbf{justificación práctica}, que sustenta la aplicabilidad operacional y el impacto en la resolución de problemas concretos; la \textbf{justificación teórica}, que establece la contribución al avance del conocimiento científico; la \textbf{justificación metodológica}, que defiende la pertinencia, validez y novedad de los métodos empleados; la \textbf{justificación ambiental}, que examina la contribución a la sostenibilidad y la reducción de impactos ecológicos; y la \textbf{justificación tecnológica}, que valora la innovación técnica y el potencial de transferencia a la práctica profesional. Referencias bibliográficas del marco de referencia: Arias Chávez, D. y Cangalaya Sevillano, L.\,M. (2023). \textit{Manual del tesista}. EDUNI, UNI; Bernal Torres, C.\,A. (2006). \textit{Metodología de la investigación} (2.\textsuperscript{a}~ed.). Pearson; Ñaupas Paitán, H. et~al. (2014). \textit{Metodología de la investigación cuantitativa-cualitativa y redacción de la tesis} (4.\textsuperscript{a}~ed.). Ediciones de la U; Vara Horna, A.\,A. (2010). \textit{¿Cómo evaluar la rigurosidad científica de las tesis doctorales?}. USMP. Si bien la normativa interna de la universidad exige la presentación de una justificación de carácter social y, en función de su relevancia, las evaluaciones de justificación analizadas a continuación permiten articular de manera más sistemática y coherente dicho requisito establecido por la normativa institucional.


{{OBSERVACIONES_AQUI}}


\section{Evaluación de la Hipótesis de Investigación}

La hipótesis de investigación es la proposición central que anticipa la relación entre variables y orienta el diseño metodológico del proyecto de tesis. En Ingeniería Civil, toda hipótesis bien formulada articula, en un primer nivel, una relación abstracta entre constructos teóricos \textit{(hipótesis conceptual)} y, en un segundo nivel, la traduce en una predicción específica, medible y verificable mediante procedimientos estandarizados \textit{(hipótesis operacional)}. La calidad de la hipótesis no depende únicamente de su correcta redacción gramatical, sino de su solidez epistemológica: debe ser clara, estar fundamentada en teoría y evidencia previa, especificar las variables con precisión operacional, ser susceptible de contrastación empírica mediante métodos estadísticos adecuados, y derivarse con coherencia lógica del problema, la pregunta y los objetivos de investigación declarados. Esta rúbrica permite evaluar, de manera sistemática e independiente de la especialidad de Ingeniería Civil involucrada (geotécnica, estructuras, pavimentos, hidráulica, materiales, mecánica de rocas, infraestructura de transporte u otras), la calidad integral de la hipótesis formulada en el proyecto de tesis de pregrado. La evaluación se estructura en cinco dimensiones complementarias que cubren la claridad enunciativa, la fundamentación teórica, la operacionalización de variables, la verificabilidad empírica y la coherencia con el planteamiento del problema. La rúbrica es aplicable tanto a hipótesis causales, correlacionales y mecanísticas, como a hipótesis experimentales unifactoriales o multifactoriales.\par


{{OBSERVACIONES_AQUI}}


\section{Evaluación de la Formulación de la Pregunta de Investigación}

La pregunta de investigación constituye el eje conceptual alrededor del cual se estructura toda la investigación empírica, orientando la selección de marcos teóricos, estrategias metodológicas y procedimientos analíticos. Su formulación representa un acto intelectual sustantivo que delimita los confines epistémicos, el alcance operacional y la viabilidad práctica de un estudio. En Ingeniería Civil, la pregunta debe traducir un problema técnico ---identificado en el planteamiento del problema--- en un enunciado interrogativo que especifique variables mensurables, poblaciones o unidades de análisis claramente definidas, condiciones de contorno y una estructura semántica compatible con el diseño metodológico previsto. Una pregunta bien formulada no es genérica ni ambigua: distingue entre preguntas descriptivas, relacionales, comparativas o explicativas, y anticipa la lógica inferencial que guiará la recolección y el análisis de datos. Esta rúbrica permite evaluar la calidad, precisión y completitud de la pregunta de investigación, verificando que sea técnicamente rigurosa, operacionalmente viable, metodológicamente coherente y científicamente pertinente como componente articulador entre la problematización, los objetivos, la hipótesis y el diseño de investigación. La evaluación se centra en la consistencia interna, la especificidad operacional y la suficiencia de la pregunta como antecedente lógico del diseño experimental o metodológico del proyecto de tesis de pregrado en Ingeniería Civil.



{{OBSERVACIONES_AQUI}}


\section{Evaluación de las Variables, Dimensiones e Indicadores}

La definición, clasificación y operacionalización de variables es un requisito metodológico en las tesis de pregrado en Ingeniería Civil, pues vincula los conceptos teóricos con las mediciones empíricas que permiten contrastar las hipótesis. Una variable es cualquier propiedad, atributo o característica que puede asumir distintos valores o categorías en un dominio de observación; su delimitación rigurosa organiza el diseño experimental, la recolección de datos y la comparación de resultados. Las dimensiones agrupan las variables en categorías conceptualmente coherentes —técnicas, económicas, ambientales, sociales, microestructurales, entre otras— que representan distintos aspectos del fenómeno estudiado, mientras que los indicadores concretan cada variable en resultados verificables mediante procedimientos normalizados. Esta rúbrica evalúa la calidad, consistencia y completitud del sistema \textbf{Variable}~$\rightarrow$~\textbf{Dimensión}~$\rightarrow$~\textbf{Indicador} del proyecto de tesis, verificando la precisión en la identificación de variables, la correcta clasificación según función y escala de medición, la coherencia y organización jerárquica de la estructura dimensional, el respaldo de los indicadores en procedimientos estandarizados de medición y la consistencia interna y suficiencia del sistema para la verificación empírica de las hipótesis. La evaluación se aplica a cualquier subárea de Ingeniería Civil—geotecnia, estructuras, pavimentos, hidráulica, materiales, gestión de la construcción, saneamiento, entre otras—y se centra en la consistencia interna, la precisión técnica y la suficiencia operacional del sistema de variables como base del diseño metodológico.


{{OBSERVACIONES_AQUI}}


\section{Revisión del Marco Teórico y Conceptual}

El marco teórico constituye la estructura conceptual que fundamenta, organiza y delimita el conocimiento existente sobre el cual se sustenta una investigación. En un proyecto de tesis de pregrado en Ingeniería Civil, el marco teórico no se puede concebir como un resumen pasivo de la literatura ni como una recopilación enciclopédica de definiciones, sino como un modelo estructurado que identifica las teorías, los principios, las variables y las relaciones que guían la formulación de hipótesis, la selección metodológica y la interpretación de resultados. Su función epistemológica es proporcionar el \textit{plano conceptual} de la investigación: así como un plano estructural dicta las cargas, los apoyos y las secciones de un elemento, el marco teórico dicta los constructos, los supuestos y los vínculos lógicos que sostienen el diseño investigativo. Esta rúbrica evalúa la calidad, coherencia y suficiencia del marco teórico de un \textbf{proyecto de tesis} de pregrado --- documento que, de ser aprobado, recién se convertirá en una tesis --- verificando que la fundamentación teórica sea pertinente al problema formulado, que las fuentes sean actuales y confiables, que las variables estén conceptualmente definidas, que exista alineación verificable entre el marco teórico y los demás componentes del proyecto, y que el nivel de profundidad sea el apropiado para una investigación de pregrado en Ingeniería Civil. La evaluación se sustenta en criterios derivados de marcos de calidad reconocidos internacionalmente (Lovitts, 2007; Boote \& Beile, 2005; Grant \& Osanloo, 2014), estándares de acreditación (ABET, EUR-ACE) y estudios específicos de evaluación de tesis en programas de Ingeniería Civil latinoamericanos (García-Ramírez, 2023; Del Savio et~al., 2024; Oyewobi et~al., 2024).



{{OBSERVACIONES_AQUI}}


\section{Evaluación del Diseño Metodológico}

La metodología constituye el componente del proyecto de tesis que operacionaliza la estrategia de investigación: traduce el problema, los objetivos y las hipótesis en un plan de acción técnicamente riguroso, replicable y verificable. En Ingeniería Civil, la metodología se distingue de la tradición de las ciencias sociales en aspectos fundamentales: las variables son magnitudes físicas medibles con instrumentos calibrados; los especímenes son materiales, probetas, muestras de suelo, elementos estructurales o tramos de infraestructura; los protocolos de ensayo están normalizados por organismos como ASTM, AASHTO, NTP, ISO o ACI; y la inferencia se sustenta en principios de mecánica de suelos, mecánica de materiales, hidráulica, análisis estructural u otras ramas reconocidas de la disciplina. Esta rúbrica evalúa la calidad, consistencia, completitud y rigor técnico del diseño metodológico propuesto, verificando que el tesista haya planificado un procedimiento experimental o analítico capaz de responder a la pregunta de investigación con validez interna, trazabilidad de datos y control de las fuentes de variabilidad propias de los materiales y procesos de la Ingeniería Civil.\par


{{OBSERVACIONES_AQUI}}


\section{Evaluación de Cientificidad del Proyecto}

La cientificidad de un proyecto de tesis de pregrado en Ingeniería Civil es la propiedad que permite distinguir un planteamiento orientado a la \textbf{generación de conocimiento científico} de uno que se limita a la \textbf{aplicación profesional rutinaria} de normas, procedimientos o tecnologías ya establecidas. Un proyecto de tesis científicamente formulado identifica un vacío de conocimiento genuino, propone hipótesis falsables, diseña una estrategia experimental o analítica con control de variables, anticipa el tratamiento cuantitativo de incertidumbres y aspira a producir resultados transferibles más allá del sitio, material o proyecto específico estudiado. La presente rúbrica operacionaliza diez principios epistemológicos de la contribución científica en ingeniería ---comprensión mecanicista, generalizabilidad, hipótesis falsables, control de variables, cuantificación de incertidumbres, transparencia metodológica, validación cruzada, innovación teórica o metodológica, ampliación del conocimiento y difusión académica--- en cinco dimensiones evaluables con veinte criterios verificables. Su propósito es proporcionar a los evaluadores un mecanismo sistemático y reproducible para determinar si un proyecto de tesis posee el nivel de cientificidad suficiente para ser aprobado como investigación de pregrado en Ingeniería Civil, diferenciándolo de un informe técnico, un expediente de diseño o un trabajo de aplicación profesional. La evaluación se aplica al proyecto \textit{antes} de su ejecución; un proyecto que supere el umbral de cientificidad podrá formularse como tesis de investigación. Este instrumento se fundamenta en los diez principios de la contribución científica en ingeniería (Arbulú, 2026), que establecen las características distintivas de la investigación académica frente a la práctica profesional.


{{OBSERVACIONES_AQUI}}


\section{Evaluación de la Reproducibilidad y Trazabilidad}

La reproducibilidad ---la capacidad de que un investigador independiente obtenga resultados consistentes siguiendo el mismo protocolo--- y la trazabilidad ---la posibilidad de rastrear cada resultado hasta los datos primarios, materiales, equipos y condiciones que lo originaron--- constituyen dos pilares inseparables de la calidad científica en Ingeniería Civil experimental. En tesis de pregrado que involucran ensayos de laboratorio, campañas de campo o modelación numérica, la reproducibilidad no se logra con la mera declaración de compromiso ni con la cita genérica de normas de ensayo: exige que el proyecto especifique, con precisión operativa, la identidad química y física de los materiales empleados, el protocolo experimental paso a paso, las condiciones ambientales y de preparación de muestras, el sistema de registro de datos primarios, la suficiencia estadística del plan muestral y la ausencia de ambigüedades que impidan la réplica por terceros. Esta rúbrica evalúa el grado en que el proyecto de tesis proporciona la información técnica necesaria y suficiente para que sus procedimientos sean reproducibles y sus resultados trazables, distinguiéndose de los instrumentos previos que evalúan el compromiso declarativo con la gestión de datos (Sección~III de Observaciones Normativas) y la anticipación genérica de un sistema de trazabilidad (Dimensión~IV de Cientificidad). Aquí se verifica el contenido técnico-operativo que hace posibles la reproducibilidad y la trazabilidad, no su declaración de intención.



{{OBSERVACIONES_AQUI}}


\section{Evaluación del Presupuesto y Cronograma}

El presupuesto y el cronograma constituyen instrumentos de gestión que determinan la viabilidad operativa, económica y temporal de un proyecto de tesis de pregrado en Ingeniería Civil. Un presupuesto correctamente formulado no se limita a una lista de gastos; refleja la comprensión del tesista sobre el alcance real de su investigación, la naturaleza de los recursos requeridos y las restricciones del contexto institucional y del mercado. De forma análoga, un cronograma riguroso evidencia la capacidad de planificar la secuencia lógica de actividades, respetar dependencias técnicas, anticipar tiempos de espera y establecer hitos de control verificables. La planificación financiera y temporal deficiente produce, con frecuencia, retrasos que comprometen la calidad de los datos, matrices experimentales incompletas, omisión de ensayos o análisis sumamente importantes, o la imposibilidad de concluir el proyecto en los plazos académicos establecidos. La evaluación conjunta de ambos componentes permite verificar si el tesista ha considerado la totalidad de los gastos relevantes ---ensayos de laboratorio, ensayos de campo, materiales, software especializado, licencias, herramientas digitales, cuentas de inteligencia artificial, logística, transporte de muestras, servicios externos, consumibles y costos indirectos---, si las estimaciones son realistas y sustentables con cotizaciones o referencias de mercado, si el cronograma es coherente con la secuencia metodológica propuesta y si la planificación global demuestra un nivel de rigurosidad y sustento técnico compatible con la formación de un profesional en Ingeniería Civil. Esta rúbrica se aplica a proyectos de tesis en cualquier especialidad de la Ingeniería Civil: geotecnia, estructuras, pavimentos, hidráulica, recursos hídricos, saneamiento, gestión de la construcción, materiales de construcción, transporte, entre otras.


{{OBSERVACIONES_AQUI}}


\section{Observaciones Técnicas de Fondo}

La revisión técnica de fondo de una propuesta de tesis en Ingeniería Civil exige que el evaluador examine, en primer lugar, la consistencia entre las normas técnicas citadas y el uso efectivo que el proyecto hace de ellas. Es frecuente que los tesistas invoquen normativa internacional o nacional ---ASTM, AASHTO, NTP, ACI, Eurocódigos, entre otras--- de manera nominal, sin verificar que las condiciones de aplicabilidad establecidas por la propia norma se cumplan en el contexto específico de la investigación. El revisor debe contrastar si los rangos de validez del ensayo o del modelo normativo son coherentes con el tipo de material, la escala de estudio y las condiciones de contorno declaradas en el proyecto. Asimismo, conviene verificar que las versiones citadas de las normas correspondan a ediciones vigentes y que las eventuales modificaciones introducidas respecto del procedimiento estándar estén explícitamente justificadas. Una incongruencia típica consiste en aplicar métodos de ensayo diseñados para un tipo de material a otro cuyas propiedades físicas o químicas difieren sustancialmente, sin discutir las implicaciones de dicha extrapolación sobre la validez de los resultados.

En segundo lugar, el evaluador debe prestar atención a la coherencia entre los principios teóricos invocados en el marco conceptual y las decisiones metodológicas adoptadas en el diseño de la investigación. Un proyecto técnicamente sólido no se limita a enunciar teorías o modelos reconocidos de la disciplina, sino que demuestra que las hipótesis de partida de dichos modelos son compatibles con las condiciones del problema investigado. El revisor debe examinar si los supuestos implícitos de los modelos constitutivos, las ecuaciones de diseño o los criterios de falla seleccionados son pertinentes al fenómeno estudiado o si, por el contrario, se emplean fuera de su dominio de validez. Esta verificación resulta especialmente relevante cuando el proyecto combina enfoques provenientes de distintas subespecialidades ---por ejemplo, mecánica de suelos con química de materiales, o hidráulica con análisis estructural---, ya que la integración de marcos teóricos heterogéneos demanda una justificación explícita de su compatibilidad conceptual y operacional.

En tercer lugar, resulta útil evaluar la trazabilidad lógica entre el planteamiento del problema, los objetivos, la hipótesis y el diseño experimental, con especial énfasis en la detección de inconsistencias entre lo que el proyecto declara investigar y lo que su metodología efectivamente permite concluir. El revisor debe identificar situaciones en las que los objetivos anticipen conclusiones causales, mientras que el diseño solo admite inferencias correlacionales, o en las que la hipótesis postule relaciones entre variables que la matriz experimental propuesta no es capaz de contrastar con la potencia estadística suficiente. De igual manera, conviene verificar que las variables independientes estén definidas con la precisión necesaria para garantizar la reproducibilidad del experimento y que los niveles o tratamientos seleccionados no sean arbitrarios, sino técnicamente fundamentados. La ausencia de grupos de control, la omisión de variables intervinientes relevantes o la confusión entre la significancia estadística y la relevancia práctica para la ingeniería constituyen debilidades frecuentes que el evaluador debe señalar con claridad.

En cuarto lugar, el revisor debe examinar si los parámetros, las magnitudes y las unidades empleadas a lo largo del proyecto son internamente consistentes y si las afirmaciones técnicas están respaldadas por evidencia verificable. Es habitual encontrar proyectos que presentan valores de propiedades mecánicas, hidráulicas o térmicas sin indicar las condiciones de ensayo bajo las cuales fueron obtenidos, o que citan resultados de la literatura sin evaluar si las condiciones experimentales de los estudios de referencia son comparables con las del contexto local. El evaluador debe verificar que los órdenes de magnitud reportados sean plausibles para el tipo de material y las condiciones declaradas, que las conversiones entre sistemas de unidades sean correctas y que los criterios de aceptación o umbrales normativos contra los cuales se compararán los resultados estén explícitamente identificados. La detección de cifras incongruentes, unidades omitidas o valores fuera de los rangos típicos reportados en la literatura especializada constituye un indicador temprano de debilidades conceptuales o de falta de dominio técnico por parte del tesista.

En quinto lugar, el evaluador debe valorar si el proyecto reconoce las limitaciones inherentes a su enfoque técnico y si anticipa las fuentes de incertidumbre que podrían afectar la interpretación de los resultados. Un proyecto que no identifica las restricciones de generalización derivadas de la escala de ensayo, de la variabilidad natural de los geomateriales, de las simplificaciones introducidas en los modelos numéricos o de las condiciones ambientales no controladas revela una comprensión superficial del problema investigado. El revisor debe distinguir entre las limitaciones declarativas —aquellas que el tesista enuncia de manera genérica sin consecuencias operativas— y las limitaciones integradas, que se reflejan en decisiones concretas del diseño metodológico, como la inclusión de réplicas suficientes, la adopción de intervalos de confianza, la propagación de incertidumbres o la validación cruzada con métodos independientes. La ausencia de esta reflexión crítica compromete la credibilidad científica de la propuesta y debe ser observada como una debilidad técnica de fondo.

Finalmente, el evaluador debe considerar que las observaciones técnicas de fondo no se agotan en la verificación de errores puntuales, sino que abarcan la evaluación integral de la solidez conceptual del proyecto como contribución al conocimiento de la Ingeniería Civil. Esto implica examinar si el proyecto distingue con claridad entre un problema de conocimiento genuino y un problema de aplicación profesional rutinaria, si la estrategia investigativa propuesta tiene el potencial de generar resultados transferibles más allá del caso particular estudiado y si el diseño metodológico refleja una comprensión mecanicista del fenómeno que trasciende la mera adecuación empírica. Las incongruencias técnicas más profundas no residen en errores de cálculo o en citas normativas incorrectas, sino en la desarticulación entre los fundamentos científicos de la disciplina y la implementación operativa del proyecto. El revisor, al formular sus observaciones, debe orientar al tesista hacia la corrección de estas debilidades estructurales con indicaciones específicas y constructivas que fortalezcan la calidad científica del trabajo propuesto.


{{OBSERVACIONES_AQUI}}


\section{Estructura Documental del Proyecto de Tesis}

La estructura documental de un proyecto de tesis de pregrado en Ingeniería Civil no constituye un requisito meramente formal ni una exigencia burocrática, sino un indicador verificable de la madurez metodológica del tesista y del grado de institucionalización de su propuesta dentro del marco normativo que la universidad ha establecido para la producción de conocimiento. Un documento estructuralmente completo, ordenado y conforme a los lineamientos institucionales facilita la evaluación académica, garantiza la trazabilidad de los componentes del proyecto, permite la verificación de la coherencia interna entre secciones y asegura que el trabajo podrá ser registrado, archivado y difundido conforme a las exigencias del sistema nacional de repositorios y grados académicos. La presente rúbrica evalúa la conformidad del proyecto de tesis con el esquema aprobado por la autoridad competente, la completitud y corrección de los elementos preliminares (carátula, metadatos, índices), la presencia, secuencia y articulación de los elementos del cuerpo (planteamiento del problema, marco teórico, hipótesis, variables, metodología, aspectos administrativos), la inclusión y calidad de los elementos finales (referencias, anexos, cronograma, presupuesto, matriz de consistencia) y la coherencia global de la presentación documental como unidad integrada. La evaluación se sustenta en la normativa institucional vigente y es aplicable a proyectos de tesis de cualquier especialidad dentro de la Ingeniería Civil y se centra en la verificación objetiva de la presencia, completitud, conformidad normativa y articulación de los componentes documentales exigidos.





{{OBSERVACIONES_AQUI}}


\section{Calidad de la Redacción Académica}

La calidad de la redacción académica en un proyecto de tesis de Ingeniería Civil es condición necesaria para la comunicación, la verificabilidad y la credibilidad científica de la propuesta. Un proyecto con redacción imprecisa, ambigua o gramaticalmente deficiente dificulta la evaluación de los dictaminantes y compromete la transmisión exacta de las ideas, procedimientos y hallazgos que constituyen el aporte del tesista. La escritura científica en Ingeniería Civil se caracteriza por la precisión terminológica —cada concepto técnico tiene un significado unívoco en la disciplina—, el rigor discursivo —las afirmaciones se sustentan en evidencia y se articulan mediante una lógica argumentativa verificable— y la corrección lingüística —el texto se ajusta a las normas ortográficas, gramaticales y de puntuación del español académico—. A estas propiedades se suman la cohesión textual —articulación fluida y lógica entre oraciones, párrafos y secciones— y la adecuación al registro académico formal —adopción de convenciones de escritura científica que distinguen el discurso técnico del lenguaje coloquial, periodístico o literario—. Esta rúbrica evalúa la calidad de la redacción académica del proyecto de tesis en cinco dimensiones complementarias que abarcan la competencia comunicativa escrita exigible a un tesista de pregrado en Ingeniería Civil. La evaluación no considera el estilo personal del autor —legítimamente variable—, sino el cumplimiento de los estándares mínimos de claridad, precisión, corrección y rigor que la comunidad académica y profesional exige para que un documento técnico sea publicable, evaluable y reproducible.


{{OBSERVACIONES_AQUI}}


\section{Formato y Presentación}

El formato y la presentación de un proyecto de tesis es un requisito de comunicación técnica que garantiza la legibilidad, la trazabilidad y la verificabilidad del documento como producto académico. En el marco del Reglamento Específico para Obtener el Grado Académico, el proyecto de tesis debe ajustarse a los esquemas aprobados, emplear la carátula institucional oficial, incluir los metadatos y el informe de similitud, además debe estar eb hoja A4, fuente Times New Roman tamaño~12, interlineado de 1,5~líneas, texto justificado en ambos márgenes, márgenes de 2,5\,cm, paginación en arábigos en el extremo superior derecho (excepto la carátula), jerarquía de hasta cinco niveles de encabezados con formato diferenciado, y referencias en norma APA. La presente rúbrica evalúa la calidad formal del documento en cinco dimensiones: estructura documental y elementos obligatorios, normas tipográficas y maquetación, presentación de tablas y figuras, tratamiento de ecuaciones y expresiones matemáticas, y uso de unidades de medida y convenciones numéricas.


{{OBSERVACIONES_AQUI}}


\section{Referencias Bibliográficas}

Las referencias bibliográficas constituyen el aparato erudito que sustenta la investigación y permiten al evaluador verificar la trazabilidad, la originalidad y la fundamentación del proyecto de tesis. En el marco del Reglamento Específico para Obtener el Grado Académico, el proyecto de tesis debe emplear un sistema de citación normalizado; las Consideraciones Generales de la Propuesta para el Contenido del Proyecto de Tesis establecen de manera explícita: \textbf{``REFERENCIAS: USAR NORMA APA''}, lo que presupone un manejo riguroso y verificable de las fuentes. La Integridad Científica y Ética de la Investigación ya evalúa la ausencia de plagio, la corresponsabilidad del asesor y el sistema de citación normalizado; la presente sección complementa aquella evaluación centrándose en la \textit{calidad técnica}, la \textit{pertinencia disciplinar}, la \textit{actualidad}, la \textit{cobertura normativa} y la \textit{corrección formal} de las referencias bibliográficas como componente esencial de un proyecto de investigación en Ingeniería Civil.



{{OBSERVACIONES_AQUI}}


\section*{Veredicto Integral Final}
\noindent {{VEREDICTO_IA}}
\vspace{0.5cm}

\newpage\section*{Referencias}
\begin{list}{}{\setlength{\itemindent}{-1.27cm}\setlength{\leftmargin}{1.27cm}}
{{REFERENCIAS_IA}}
\end{list}

Atentamente,\\[1cm]
{{EVALUADOR}}

\vfill\hrule\vspace{0.2cm}\begin{center}
\textbf{PUNTAJE GLOBAL OBTENIDO: {{PUNTAJE_TOTAL_IA}}}
\end{center}\end{document}